\section{NLI as a shared backbone}\label{section:disc-nli-as-shared-backbone}

The grounding filter and the relation classifier share a single NLI model, \texttt{PubMedBERT-MNLI-MedNLI}~\parencite{deka2023pubmedbertextended}, without additional training. They differ only in their inputs: grounding uses a claim's text and its cited evidence, whereas the relation classifier compares two claims and labels them as \texttt{SUPPORT} / \texttt{CONTRADICT} / \texttt{UNRELATED} (§\ref{section:results-nli-in-grounding-and-relation-classification}). The NLI model in both roles applies fixed $0.5$ thresholds (not tuned per task), which perform acceptably across the two tasks, justifying the use of a single model.

\begin{itemize}[label={},leftmargin=1.5em]

    \item \textbf{The relation classifier performs well on the synthetic dataset, with an important caveat.} The primary \texttt{predicate\_only} mode reaches $0.927$ accuracy and $0.927$ macro-F1. Especially, \texttt{SUPPORT} predictions are reliable with a $1.000$ precision, while most errors occur at the boundary between \texttt{UNRELATED} and \texttt{CONTRADICT} pairs, as shown in Table~\ref{tab:results-relation-classification-detailed}. The entailment and contradiction thresholds were frozen at $0.5/0.5$, and were never tuned on this set; a joint diagnostic sweep over a grid of threshold values (each in $\{0.4, \dots, 0.9\}$) was run, confirming that the frozen values are within a fraction of a percentage point of the optimum (entailment threshold~0.4,
    contradiction threshold~0.9: macro-F1 0.9307). This shows that the classifier is robust to changes, and since there was no tuning, the evaluation is completely out-of-sample. 
    
    The caveat, also stated in the results, is that the synthetic pairs might consist of clearly supporting, contradicting, or unrelated examples; therefore, the obtained result should be treated as an optimistic ceiling rather than a direct estimate of production performance. Claim pairs emitted by the production pipeline are likely more difficult, especially when two findings appear similar on the surface but are about different diseases or entities. 
    
    \item \textbf{Scope-aware input is a design choice that the evaluation cannot test.} The \texttt{scope\_aware} mode, which prepends a disease/entity scope to each claim for context, scored marginally worse on the synthetic dataset (\texttt{predicate\_only}'s $0.927$ versus \texttt{scope\_aware}'s $0.923$), which confirms that the bracketed scope prefix does not degrade the classifier. However, the synthetic set cannot test the mechanism, as each claim pair shares the same disease/entity label, making the scope prefix of the claims identical. The real corpus, where rules carry their own scope information and can differ, emits too few relations to test and compare the real-world performance of the two different modes (§\ref{section:disc-knowledge-extraction-quality-and-stability}). Therefore, scope-awareness is untested, and validating it on production rules is left to future work.

\end{itemize}
    
The broader conclusion is that an off-the-shelf biomedical NLI model with reasonable thresholds performs sufficiently well, both as a faithfulness guard and a relation classifier. The reliability of the relation classifier on the more ambiguous cases, however, remains untested for the reason above.